% !TeX spellcheck = de_DE

Zur Untersuchung des in Kapitel \ref{ch:konzeption} beschriebenen Ansatzes wurde ein Proof of Concept (PoC) entwickelt, der einen auf Metriken basierten Feedbackloop für Texte in Leichter Sprache Plus (LS+) realisiert.
Ziel der prototypischen Umsetzung ist nicht die Entwicklung eines produktionsreifen Systems, sondern die Bereitstellung einer technischen Basis, die zur Evaluation der Forschungsfragen herangezogen werden kann.

Der Prototyp ermöglicht die automatisierte Bewertung generierter LS+-Texte anhand definierter Metriken sowie deren iterative Überarbeitung mittels feedbackbasiertem Re-Prompting und einem LLM.
Dadurch kann untersucht werden, ob dieses Vorgehen zu einer messbaren Verbesserung der Textqualität führt.

Folgendes Kapitel wird die Umsetzung des Prototyps beleuchten und die Designentscheidungen aufgreifen um einen Überblick über die technische Grundlage zu geben.